\section{โครงสร้างและวิธีออกแบบ}
\subsection{สถาปัตยกรรมข้อมูล}
โค้ดมีโครงสร้าง \nolinkurl{wh_*} เดิมใน Prisma และโครงสร้าง \nolinkurl{dw_*} รุ่นใหม่ในชุด ETL อยู่ร่วมกัน บทความนี้ใช้รุ่นใหม่เป็นแกน เพราะแยกใบลาออกจากรายวิชาและรองรับบริบทหน่วยงานชัดเจน การนับแถวในตารางเดิมจึงไม่ควรนำมาเทียบกับจำนวนใบลาในรุ่นใหม่โดยไม่ตรวจ grain\cite{model,etl}

ชุด ETL อยู่ใน \nolinkurl{apps/api/scripts/warehouse-etl} โดยแยกโมดูลสกัด แปลง ตรวจสอบ และโหลดข้อมูล บริการ \nolinkurl{WarehouseLiveService} อ่านตารางรุ่นใหม่และจัดขอบเขตตามผู้ใช้ ส่วน \nolinkurl{live-report-data.ts} แปลงชุดข้อมูลเป็นโครงสร้างที่หน้ารายงานใช้งานได้

\subsection{ตารางข้อเท็จจริง}
\textbf{ระดับใบลา:} \nolinkurl{dw_fact_leave_request} กำหนดหนึ่งแถวต่อหนึ่งใบลาจากระบบต้นทาง มีรหัสต้นทางเพื่อป้องกันรายการซ้ำ เชื่อมนิสิต ภาคการศึกษา ประเภทลา หน่วยงาน หลักสูตร และที่ปรึกษา เก็บวันเริ่ม วันสิ้นสุด จำนวนวันลา ชั้นปี และข้อมูลกำกับรอบนำเข้า ยอดรวมระดับนิสิต ภาควิชา หรือคณะควรเริ่มจากตารางนี้

\textbf{ระดับใบลากับกลุ่มเรียน:} \nolinkurl{dw_fact_leave_subject} มีคีย์ร่วมระหว่างใบลากับ offering หนึ่งใบลาจึงมีได้หลายแถว เก็บจำนวนวันเรียน ครั้งเรียน และชั่วโมงเรียนที่ได้รับผลกระทบเมื่อมีข้อมูลเพียงพอ ใช้วิเคราะห์รายวิชาและกลุ่มเรียน ไม่ใช้จำนวนแถวแทนจำนวนใบลาของทั้งคณะ

ตัวอย่างเชิงอธิบาย: ใบลา A มี 2 วันและเกี่ยวข้องกับ 3 รายวิชา ยอดภาพรวมเท่ากับ 1 ใบและ 2 วัน ส่วนความสัมพันธ์รายวิชาเท่ากับ 3 รายการ หากผู้สอนร่วมมี 2 คน จำนวนแถวหลังเชื่อมอาจเพิ่มอีก การรวมต้องกลับไปที่คีย์ใบลาหรือคีย์ใบลา--offering ตามคำถาม ตัวอย่างนี้มิใช่สถิติจากฐานข้อมูลจริง

\subsection{มิติและตารางเชื่อม}
มิตินิสิตและบุคลากรระบุตัวตน มิติประเภทลาจัดกลุ่มเหตุการณ์ มิติภาคการศึกษากำหนดลำดับเทอมและปีสำหรับเปรียบเทียบ แม้เทอมนั้นไม่มีใบลา มิติรายวิชาแยกรหัสและรุ่นข้อมูลวิชา ส่วนมิติ offering ระบุรายวิชา เทอม กลุ่มเรียน และชนิดกลุ่มเรียน

มิติคณะ ภาควิชา และหลักสูตรอธิบายโครงสร้างองค์กร โดยสังกัดนิสิตใน fact แยกจากเจ้าของรายวิชาในมิติ subject/offering และสังกัดอาจารย์ ตาราง \nolinkurl{dw_snapshot_student_term} รองรับบริบทนิสิตรายเทอม ส่วนตารางเชื่อมบุคลากรกับหน่วยงานรายเทอมรองรับการจัดกลุ่มบุคลากร

ตาราง \nolinkurl{dw_bridge_offering_lecturer} รองรับผู้สอนร่วมของกลุ่มเรียน และ \nolinkurl{dw_bridge_leave_subject_lecturer} เก็บความสัมพันธ์ผู้สอนกับรายวิชาในใบลา การกรองรายงานของผู้สอนจึงต้องระวังไม่ทำให้ใบลาเดียวถูกนับซ้ำตามจำนวนผู้สอน

\subsection{กระบวนการ ETL และคุณภาพข้อมูล}
การสกัดรวบรวมรายการต้นทางและข้อมูลอ้างอิง การแปลงจัดรูปแบบรหัส วันที่ ประเภทลา และความสัมพันธ์ให้เป็นแบบเดียวกัน การโหลดแก้ไขหรือเพิ่มข้อมูลด้วยคีย์ต้นทาง และบันทึกรอบนำเข้าไว้ใน \nolinkurl{dw_load_batch} ส่วน \nolinkurl{dw_data_issue} รองรับการติดตามปัญหาคุณภาพข้อมูล

กระบวนการปัจจุบันเป็นการประมวลผล snapshot ทั้งชุดในหน่วยความจำ แล้วเขียนเป้าหมายใน transaction เดียว การประกาศรอบสำเร็จจึงผูกกับข้อมูลที่เขียนครบ ไม่ใช่กระบวนการ CDC หรือการส่งข้อมูลทันทีทุกเหตุการณ์ การจัดการรายการที่หายจากต้นทางต้องใช้ reconciliation ตามเงื่อนไขที่กำหนด

รหัสนิสิต รายวิชา และรหัสต้นทางต้องคงรูปสตริงเพื่อรักษาเลขศูนย์นำหน้าและหลีกเลี่ยงเลขขนาดใหญ่สูญเสียความแม่นยำ การจับคู่สังกัด หลักสูตร และบุคลากรต้องมีข้อมูลอ้างอิงที่ตรวจสอบได้ ค่าที่ไม่ทราบต้องแยกจากศูนย์ โดยเฉพาะผลกระทบต่อการเรียนที่ยังไม่มีตารางเรียนครบถ้วน

\subsection{นิยามตัวชี้วัด}
จำนวนใบลาคือจำนวนคีย์ใบลาที่ไม่ซ้ำในขอบเขตที่เลือก จำนวนผู้ลาคือจำนวนนิสิตไม่ซ้ำ ส่วนวันลารวมเป็นผลรวม \nolinkurl{request_days} จาก fact ระดับใบลา การใช้ผลรวมค่าที่ไม่ซ้ำตามตัวเลขวันลาไม่ถูกต้อง เพราะหลายใบอาจมีจำนวนวันเท่ากัน

วันลาตามปฏิทินไม่เท่ากับวันเรียนหรือชั่วโมงเรียนที่ขาด การวัดผลกระทบรายวิชาต้องอาศัยวันที่เรียนจริง ระยะเวลาคาบ และความสัมพันธ์กับใบลา การรวมชั่วโมงหลายวิชามีหน่วยเป็นชั่วโมงเรียนที่ได้รับผลกระทบ ไม่ใช่เวลาตามปฏิทินของนิสิต

ค่าเฉลี่ยต่อเทอมต้องใช้จำนวนเทอมที่เลือกเป็นตัวหาร รวมเทอมยอดศูนย์ สำหรับ box plot ของรายวิชา หน่วยตัวอย่างเป็นยอดรวมของรายวิชาต่อเทอมในช่วงที่เลือก มิใช่ความยาวของใบลาแต่ละใบ และต้องแสดงกรณีไม่มีตัวอย่างอย่างชัดเจน

\subsection{ขอบเขตการเข้าถึง}
บริการรายงานตรวจตัวตนและบทบาทจากฝั่งเซิร์ฟเวอร์ แล้วใช้รายชื่อนิสิตในที่ปรึกษา รายวิชาและกลุ่มเรียนที่สอน หรือหน่วยงานปัจจุบันเป็นขอบเขต ผู้บริหารเข้าถึงหน่วยงานภายในคณะที่ได้รับสิทธิ์ ส่วนเจ้าหน้าที่ใช้ขอบเขตภาควิชาจากบัญชีในโค้ดที่ตรวจ ไม่อนุมานว่ามีสิทธิ์ทุกคณะ การเปิดรายละเอียดต้องรักษาขอบเขตเดียวกับหน้าสรุป\cite{service}

ภาพสคีมาจากบทความต้นฉบับแสดงในภาคผนวก ข โดยภาพที่~\ref{fig:warehouse-schema-overview} แสดงการจัดกลุ่มตาราง และภาพรายละเอียดแสดงคีย์กับคอลัมน์สำหรับอ่านประกอบแบบจำลอง
